\documentclass[10pt,a4paper]{amsart}
\usepackage[margin=19mm]{geometry}
\usepackage{amsmath,amssymb,mathtools}
\usepackage[scr=rsfs,cal=euler]{mathalfa}
\usepackage{xfrac}
\usepackage[unicode,hidelinks,bookmarks=false]{hyperref}
\newcommand{\ie}{i.e.\ }
\newcommand{\cf}{cf.\ }
\newcommand{\resp}{respectively\ }
\newcommand{\Subsection}{\subsection}
\providecommand{\nobreakdash}{\nobreak\hbox{-}}
\setlength{\emergencystretch}{2em}
\multlinegap0pt
\usepackage{bm}
\usepackage{bbm}
\usepackage{dsfont}
\usepackage{xspace}
\usepackage{cancel}
\usepackage{mathtools}
\usepackage{amsthm}
\makeatletter
\@ifundefined{theorem}{\newtheorem{theorem}{Theorem}}{}
\@ifundefined{assumption}{\newtheorem{assumption}[theorem]{Assumption}}{}
\@ifundefined{lemma}{\newtheorem{lemma}[theorem]{Lemma}}{}
\makeatother
\title[Convex Distance Operator Transport: A Convex and Geometry-Pr]{Convex Distance Operator Transport: A Convex and Geometry-Preserving Formulation}
\author{Junhyoung Chung, Euijong Song, Won Hwa Kim, Gunwoong Park}
\date{Source: arXiv:2606.02047v1 (CC BY 4.0)}
\begin{document}
\maketitle

\noindent\emph{Source and licence.} Convex Distance Operator Transport: A Convex and Geometry-Preserving Formulation. Version arXiv:2606.02047v1 (CC BY 4.0), distributed under the Creative Commons Attribution 4.0 International licence, as recorded in the arXiv licence metadata for this exact version and shown on the abstract page https://arxiv.org/abs/2606.02047v1. Full rights notice: licenses/arxiv-2606.02047-CC-BY-4.0.txt.\par\smallskip

\noindent\textit{Editorial scope.} Counted unit: the Lemma whose source label is \texttt{lem:laguerre} in the article (source0.tex lines 1323--1344, source label \texttt{lem:laguerre}), delivered together with the article's own complete original proof of it (source0.tex lines 1345--1405, one \texttt{proof} environment of 61 lines). No mathematics is written, repaired or re-derived: statement and proof are the article's own text line for line. Required context retained uncounted: ass:non-atomic (lines 631--634); ass:no-tie (lines 635--642). Every definition, display and label of those lines is retained unchanged. Omitted from this package and not redistributed: 1--630, 643--1322, 1406--2443. Nothing inside a retained excerpt is omitted or altered: every retained body is delivered exactly as the article's lines are written, so no omission receipt of any kind accompanies this package. Citations: the retained excerpts carry no citation command at all, so no bibliography entry of the article is transcribed and no reference is redistributed. Reference adaptation: none was needed and none was made; every reference inside the delivered unit resolves inside it, so no unresolved reference or citation is produced. Printed slips kept literal: no printed word, symbol, hypothesis or number of the counted statement or of its proof was altered. Layout and class: the article's own class is \texttt{article} with packages including microtype, graphicx, subcaption, booktabs, multirow, array, pifont, float, hyperref, icml2026, amsmath, amssymb, mathtools, amsthm; the wrapper is the lane's pinned \texttt{amsart} editorial wrapper at 10pt on A4, which restates the theorem-like environments the retained text needs and adds the article's own macro definitions that the retained text needs (none), with extra wrapper packages (bm, bbm, dsfont, xspace, cancel, mathtools). The delivered numbering is the unit's own. Assets: no retained span contains a figure environment, a graphics-inclusion command, a PDF-page inclusion, a table or a TikZ picture; no image file is copied into this package, the package declares no asset dependency and no image file is redistributed with it. Licence: version arXiv:2606.02047v1 of this article is distributed under the Creative Commons Attribution 4.0 International licence, as recorded in the arXiv licence metadata for this exact version and shown on the abstract page https://arxiv.org/abs/2606.02047v1. A full rights notice is delivered at licenses/arxiv-2606.02047-CC-BY-4.0.txt. Characters: every retained excerpt is pure ASCII, so the delivered engine, XeLaTeX with its default Latin Modern text font, reports no missing character.
\begin{assumption}[No atoms]
	\label{ass:non-atomic}
	$\mathbb{P}_X$ and $\mathbb{P}_Y$ are non-atomic; $\mathbb{P}_X(\{x\}) = \mathbb{P}_Y(\{y\}) = 0$ for all $x \in \mathcal{X}$ and $y \in \mathcal{Y}$.
\end{assumption}

\begin{assumption}[No ties]
	\label{ass:no-tie}
	For any distinct $x_1,x_2 \in \mathcal{X}$, $y_1,y_2 \in \mathcal{Y}$, and for any $c \in \mathbb{R}$,
	\begin{align*}
		&\mathbb{P}_X(\{x \in \mathcal{X}: d_{\mathcal{X}}(x,x_1) - d_{\mathcal{X}}(x,x_2) = c\}) = 0 ,\\
		&\mathbb{P}_Y(\{y \in \mathcal{Y}: d_{\mathcal{Y}}(y,y_1) - d_{\mathcal{Y}}(y,y_2) = c\}) = 0 .
	\end{align*}
\end{assumption}

\begin{lemma}[Laguerre lemma]
	\label{lem:laguerre}
	Let $(\mathcal{X},d_{\mathcal{X}})$ be a compact metric space, $\mathbb{P}_X$ be a Borel probability measure on $\mathcal{X}$ that satisfies Assumptions~\ref{ass:non-atomic} and \ref{ass:no-tie}, and $\hat{\mathbb{P}}_X$ be the empirical probability measure of $\mathbb{P}_X$ supported on $n$ fixed distinct data points, i.e.,
	\begin{align*}
		\hat{\mathbb{P}}_X = \frac{1}{n} \sum_{i=1}^{n} \delta_{X_i} .
	\end{align*}
	Let $\hat{\mathcal{X}} \coloneqq \{X_1,...,X_n\} \subset \mathcal{X}$. Then $(\hat{\mathcal{X}},\hat{\mathbb{P}}_X)$ is clearly a Polish probability space. Let $(\pi^*,\psi^*)$ be primal and dual optimal solutions for the semi-discrete problem:
	\begin{align*}
		\min_{\pi \in \Pi(\mathbb{P}_X,\hat{\mathbb{P}}_X)}\left(\int_{\mathcal{X} \times \hat{\mathcal{X}}} d_{\mathcal{X}}(x,\hat{x}) \pi(dx,d\hat{x})\right) = \max_{\psi \in L^1(\hat{\mathcal{X}},\hat{\mathbb{P}}_X)} \left( \int_{\mathcal{X}} \psi^{d_{\mathcal{X}}}(x) \mathbb{P}_X(dx) - \frac{1}{n}\sum_{i=1}^{n} \psi(X_i) \right) ,
	\end{align*}
	where the ${d_{\mathcal{X}}}$-transform is taken as $\psi^{d_{\mathcal{X}}}(x) = \min_{1 \le j \le n} (d_{\mathcal{X}}(x,X_j) + \psi(X_j))$. Note that all elements in $L^1(\hat{\mathcal{X}},\hat{\mathbb{P}}_X)$ can be identified with $n$-dimensional real-valued vectors.
	
	For each $i= 1,...,n$, define the Laguerre-like cell $V_i$ by
	\begin{align*}
		V_i \coloneqq \left\{x \in \mathcal{X} : d_{\mathcal{X}}(x,X_i) + \psi^*(X_i) \leq d_{\mathcal{X}}(x,X_j) + \psi^*(X_j) , \, \forall j\right\} \cup \{X_i\} .	
	\end{align*}
	Then $\{V_i\}_{i=1}^n$ forms a $\mathbb{P}_X$-a.e. partition of $\mathcal{X}$, satisfies $\mathbb{P}_X(V_i)=1/n$ for all $i$, and
	\begin{align*}
		\pi^* = \sum_{i=1}^{n} \mathbb{P}_X\lfloor_{V_i} \otimes \delta_{X_i} , \quad \mathbb{P}_X\lfloor_{V_i}(dx) = \mathbb{I}_{V_i}(x)\mathbb{P}_X(dx) ,
	\end{align*}
	where $\mathbb{I}_{V_i}$ is the indicator function of $V_i$.
\end{lemma}

\begin{proof}
	Denote by $\psi^{*}_j \coloneqq \psi^*(X_j)$. Then, we have the following:
	\begin{align*}
		(\psi^*)^{d_{\mathcal{X}}}(x) = \min_{1 \leq j \leq n} \bigl(d_{\mathcal{X}}(x,X_j) + \psi^*_j\bigr).		
	\end{align*}
	
	Now define, for each $i=1,...,n$,	
	\begin{align*}		
		\tilde{V}_i \coloneqq \left\{x \in \mathcal{X} : d_{\mathcal{X}}(x,X_i) + \psi^{*}_i \leq d_{\mathcal{X}}(x,X_j) + \psi^{*}_j , \, \forall j\right\}.	
	\end{align*}	
	By the no-ties assumption (Assumption~\ref{ass:no-tie}), the sets $\{\tilde{V}_i\}_{i=1}^n$ form a $\mathbb{P}_X$-a.e. partition of $\mathcal{X}$, and moreover for $\mathbb{P}_X$-a.e. $x$ there exists a unique index $i(x)$ such that $x \in \tilde{V}_{i(x)}$.
	
	Let $\pi^* \in \Pi(\mathbb{P}_X,\hat{\mathbb{P}}_X)$ be a primal optimal plan. Since the second marginal of $\pi^*$ is supported on the finite set $\{X_1,...,X_n\}$, there exist Borel-measurable functions $p_i : \mathcal{X} \to [0,1]$ such that	
	\begin{align*}		
		\sum_{i=1}^n p_i(x)=1 \quad \text{for }\mathbb{P}_X\text{-a.e. }x,		
	\end{align*}	
	and the disintegration of $\pi^*$ with respect to the first marginal $\mathbb{P}_X$ takes the form	
	\begin{align*}		
		\pi^*(dx,d\hat{x}) = \sum_{i=1}^n \underbrace{p_i(x) \, \delta_{X_i}(d\hat{x})}_{= \pi^*(d\hat{x} \mid x)} \, \mathbb{P}_X(dx) .		
	\end{align*}
	
	We next use complementary slackness. By the Kantorovich duality lemma, the following holds:	
	\begin{align*}		
		(\psi^*)^{d_{\mathcal{X}}}(x) - \psi^*(\hat{x}) = d_{\mathcal{X}}(x,\hat{x}) , \quad\text{for }\pi^*\text{-a.e. }(x,\hat{x}).		
	\end{align*}	
	In particular, for each $i$ and for $\mathbb{P}_X$-a.e. $x$ such that $p_i(x)>0$, we must have	
	\begin{align*}		
		(\psi^*)^{d_{\mathcal{X}}}(x) - \psi^*(X_i) = d_{\mathcal{X}}(x,X_i).		
	\end{align*}	
	Equivalently,	
	\begin{align*}	
		(\psi^*)^{d_{\mathcal{X}}}(x) = d_{\mathcal{X}}(x,X_i) + \psi_i^*.		
	\end{align*}	
	But we already have	
	\begin{align*}		
		(\psi^*)^{d_{\mathcal{X}}}(x)=\min_{1 \leq j \leq n}\bigl(d_{\mathcal{X}}(x,X_j)+\psi_j^*\bigr),		
	\end{align*}	
	so the equality implies that $i$ attains the minimum, i.e.,	
	\begin{align*}		
		p_i(x)>0 \quad\Longrightarrow\quad x\in \tilde{V}_i \quad \text{for }\mathbb{P}_X\text{-a.e. }x.		
	\end{align*}	
	By the no-ties assumption (Assumption~\ref{ass:no-tie}), the minimizer is unique for $\mathbb{P}_X$-a.e. $x$, hence for $\mathbb{P}_X$-a.e. $x$ there exists a unique index $i(x)$ such that	
	\begin{align*}		
		p_{i(x)}(x)=1,\quad p_j(x)=0 , \quad (j\neq i(x)),		
	\end{align*}	
	and necessarily $x\in \tilde{V}_{i(x)}$. Therefore,	
	\begin{align*}		
		\pi^*(dx,d\hat{x}) = \sum_{i=1}^n \mathbb{I}_{\tilde{V}_i}(x)\,\mathbb{P}_X(dx)\,\delta_{X_i}(d\hat{x}) = \sum_{i=1}^n \mathbb{P}_X\lfloor_ {\tilde{V}_i}(dx)\,\delta_{X_i}(d\hat{x}) .		
	\end{align*}
	
	Now we prove the mass constraint $\mathbb{P}_X(\tilde{V}_i)=1/n$. Since the second marginal of $\pi^*$ is $\hat{\mathbb{P}}_X$, we have for each $i$,	
	\begin{align*}		
		\frac{1}{n} = \hat{\mathbb{P}}_X(\{X_i\}) = \pi^*(\mathcal{X}\times\{X_i\}) = \pi^*(\tilde{V}_i\times\{X_i\}) = \mathbb{P}_X(\tilde{V}_i) .		
	\end{align*}
	
	Finally, define	
	\begin{align*}		
		V_i \coloneqq \tilde{V}_i \cup \{X_i\}.		
	\end{align*}	
	Since $\mathbb{P}_X$ is non-atomic, $\mathbb{P}_X(\{X_i\})=0$, hence $\mathbb{P}_X(V_i)=\mathbb{P}_X(\tilde{V}_i)=1/n$. Moreover, the representation of $\pi^*$ remains valid when $\tilde{V}_i$ is replaced by $V_i$, because $\mathbb{P}_X \lfloor_{V_i}$ and $\mathbb{P}_X \lfloor_{\tilde{V}_i}$ are equivalent as they differ only on a $\mathbb{P}_X$-null set. This completes the proof.
\end{proof}

\end{document}
